\section{Итог спуска: разделение целых чисел ядра и непрерывного языка макромира}
\label{sec:coords-bridge}

Подведём итог пройденного пути.
Спускаясь от макроскопических законов сохранения и светового конуса через квантовую механику и предел Бекенштейна, мы обнаружили, что на самом нижнем этаже исчезают не физические законы, а лишь непрерывные вещественные числа, которыми их записывали на больших масштабах.
В ядре микроуровня~$\layerMicro$ работают только конечные целочисленные регистры $(\symMu, \phaseVarphi_{\mathrm{disc}})$ и шаги $(dh, dt)$, тогда как трансцендентные константы $\symPi$ и $\ln 2$ вместе с размерными эталонами СИ возникают лишь при переводе этого счёта на язык макроскопического прибора~$\layerMacro$ (\cref{tab:m-t-bridge}).

\begin{table}[ht]
\centering
\caption{На каком слое живёт каждая величина, в каких единицах она измеряется и на какой границе происходит перевод.}
\label{tab:m-t-bridge}
\footnotesize
\setlength{\tabcolsep}{4pt}
\begin{tabular}{@{}p{0.22\linewidth}p{0.25\linewidth}p{0.24\linewidth}p{0.23\linewidth}@{}}
\toprule
Величина & Ядро $\layerMicro$ (целые $\spaceZ$) & Граница декодера $\layerMicro\leftrightarrow\layerMacro$ & Макроуровень $\layerMacro$ / СИ \\
\midrule
пространство и время & шаги $N\in\spaceZ$ ($dh$), такты $\symTau\in\spaceZ$ ($dt$); $c_0 = 1\,dh/dt$ & $1\,dh \mapsto \lengthPlanck$, $1\,dt \mapsto \htick = \caKappaGeom\timePlanck$ & $r = N\lengthPlanck$ [м], $t = \symTau\htick$ [с], $\speedC = \caKappaGeom\caSignalSpeed$ [м/с] \\
амплитуда спинора & $\symMu\in\{1,\ldots,64\}$ в квантах амплитуды (дно $1$, потолок $K_P=64$) & деление на $K_P=64$ только в декодере: $|z|=\symMu/64$ & $|z|\in[1/64,\,1]$; плотность $\rho_E = |z|^2 u_P$ [Дж/м$^3$] \\
квадрат амплитуды и голономия & $R^2=\symMu^2$ и $\zeta_{\mathrm{real}},\zeta_{\mathrm{imag}}\in\spaceZ$ в $\mathrm{кванты}^2$ ($1\ldots 4096$) & внутри ворот отношение $\frac{\zeta_{\mathrm{imag}}}{\zeta_{\mathrm{real}}+R^2+1}$ безразмерно & $\rho = R^2/4096 \in [2^{-12},1]$ \\
плотность $\rho$ (\Axref{ax:A7}) & $\rho_q = \symMu_0^2 + \symMu_1^2 \in \spaceZ$ в $\mathrm{кванты}^2$; легальное множество --- четверть-диск $\rho_q \le K_P^2 = 4096$ в $\{1,\ldots,64\}^2$; внутри шага~$g$ клампа нет, так как шаг~$\symMu$ не трогает & проекция \texttt{scale\_a7} на четверть-диск $\rho_q \le 4096$ с симметричным округлением $+0{,}5$ на мосту $\layerMacro\to\layerMicro$ (биекция на четверть-диске) & $SU(2)$-инвариантная спинорная плотность $\rho = |z_0|^2 + |z_1|^2 = \rho_q/4096 \le 1$ \\
перекрытие / исключение (\Axref{ax:A16}) & в полярной схеме $[\symMu_0,\phaseVarphi_0,\symMu_1,\phaseVarphi_1]$ перекрытие дорожек --- инвариант движения; \Axref{ax:A16} реализована как фазовый импульс $\symCapitalPhi_{\mathrm{Pauli}} = \phaseRingCount/2 = 256$ тиков при $\symMu_0^2,\symMu_1^2 > K_P^2/2$ и $|\phaseVarphi_1-\phaseVarphi_0| < 41$ и проверяется фазой & перевод $\symCapitalPhi_{\mathrm{Pauli}} = 256$ тиков в полуоборот $\symPi$ рад при чтении траектории & фазовое отталкивание Паули $\Delta\phaseVarphi = \symPi$ ($e^{i\symPi}=-1$) \\
фаза и квант действия & $\phaseVarphi_{\mathrm{disc}},\symCapitalPhi\in\spaceZ_{\phaseRingCount}$ в тиках кольца ($s_0=1$ тик; полный цикл $h = \phaseRingCount s_0 = 512\,s_0$ без $2\symPi$) & деление полного оборота на длину окружности $2\symPi$ только при переходе к радианам: $\planckHbar := h/(2\symPi)$ & угол $\phaseVarphi\in(-\symPi,\symPi]$ [рад]; $\planckHbar = \frac{512\,s_0}{2\symPi}\approx 81{,}49\,s_0$ [Дж$\cdot$с] \\
порог Гейзенберга и кванты обмена & $\symCapitalDelta\symPhi_{\mathrm{disc}} = \frac{512+13+8}{13} = 41$ тик ($41\,s_0$); $n_E = \lfloor|\symCapitalPhi|/41\rfloor$ & $41\text{ тик} \mapsto \symCapitalDelta\symPhi_{\min} = \tfrac{1}{2}$ рад $\mapsto \ladderAction = \planckHbar/2$ & $E = n_E\ladderEnergy$ [Дж], $p = n_p\ladderMomentum$, $F = n_F\ladderForce$ \\
разрядность ячейки & слово $\mathrm{mod\_bits} = 1+|\triangle| = d^2 = 9$ бит ($\phaseRingCount=512$) & сверка с термодинамической энтропией на выходе & проверка Бекенштейна: $9 \le B_{\hv} = \frac{2\symPi}{\ln 2}\approx 9{,}065$ \\
\bottomrule
\end{tabular}
\end{table}

Из табл.~\ref{tab:m-t-bridge} следуют два жёстких правила разграничения единиц:
\begin{enumerate}[nosep]
\item \textbf{Внутри шага~$g$ на~$\layerMicro$ нет вещественных дробей и делений на~$64$:} элементарный квант амплитуды на микроуровне равен целому числу~$1$ (а не $1/64$), поэтому $R^2$, $\zeta_{\mathrm{real}}$, $\zeta_{\mathrm{imag}}$ и единичный регуляризатор знаменателя~$+1$ измеряются в одних и тех же квадратах квантов, а множитель $K_P = 64$ тика стоит только перед их безразмерным отношением, переводя его в тики фазового кольца~$\spaceZ_{\phaseRingCount}$.
\item \textbf{Перевод $\div 64$ и $\times 2\symPi/512$ выполняется ровно один раз} --- на границе декодирования состояния при выходе к макроописанию~$\layerMacro$. Смешение целых квантов ядра с нормированными к потолку Планка долями $|z|\in[0,1]$ внутри самого шага~$g$ занижает голономию вакуумного дна в $64$ раза и завышает порог знаменателя в $4096$ раз.
\end{enumerate}

На этом спуск от приборной физики к дискретному носителю завершён.
В следующих главах мы начнём обратный подъём: выберем единственную трёхмерную решётку~$\carrierLattice$, удовлетворяющую требованиям изотропии и плотнейшей упаковки (\latGCKshort), сформулируем абсолютные условия~$A_1$--$A_{16}$ и запишем явный обратимый закон эволюции~$g$.
